(a): in Gauge Theory Gravity in $d$ dimensions the Yang-Mills density is $\det(\hb)^{-1}\langle\mathbf{F}^{2}\rangle_{0}$, with $\mathbf{F} = \hb^{T}(F)$ the physical field strength of the background bivector field $F = \tfrac12 F_{\mu\nu}\,e^{\mu}\wedge e^{\nu}$ (in components, a labelled reading: $\sqrt{|g|}\,g^{\mu\alpha}g^{\nu\beta}F_{\mu\nu}F_{\alpha\beta}$, the integrand of $F \wedge \star F$). Under $g \mapsto \lambda^{2}g$, that is $\hb^{-1} \mapsto \lambda\,\hb^{-1}$ and $\hb^{T} \mapsto \lambda^{-1}\hb^{T}$, the volume factor $\det(\hb)^{-1}$ scales as $\lambda^{d}$ and the outermorphism $\hb^{T}$ on bivectors as $\lambda^{-2}$, so $\langle\mathbf{F}^{2}\rangle_{0}$ scales as $\lambda^{-4}$, the action $\int\det(\hb)^{-1}\langle\mathbf{F}^{2}\rangle_{0}\,d^{d}x$ as $\lambda^{d-4}$, and $1/g^{2}$ carries the compensating dimension. (b) is a statement about the record the series takes its couplings from. (c) rests on Theorem \claimref{dark-matter.residue-soldered-invariant-not-algebraic-one} for the inequivalence, on \claimref{dark-matter.p2-indifference}, which equates sector labels along the reversal orbit only, on \claimref{synthesis.common-ceiling-coupling-premise}, which equates the ceiling couplings along that orbit only, and on Theorem \claimref{dark-matter.mirror-tower-copy-one}(iv) for the reversal orbit. The remark at the technical level, that any rule would be a further premise, is a reading of (a) to (c) and is not proved here.
